\textbf{结论名称：}含根式的最值问题（平方消根/均值不等式/三角换元）

\textbf{适用条件：}函数式中含有二次根号，且根号下为一次或二次式，定义域为有限闭区间；常用方法包括平方消根、均值不等式或三角代换。

\textbf{变量说明：}设根号内非负表达式为 $u(x),v(x)\ge0$；若用三角换元，常令 $x=a\sin\theta$ 处理 $\sqrt{a^{2}-x^{2}}$，令 $x=a\tan\theta$ 处理 $\sqrt{x^{2}+a^{2}}$ 等。

\textbf{核心公式：}
\[
\boxed{(\sqrt{u}+\sqrt{v})^{2}=u+v+2\sqrt{uv}\le 2(u+v),\qquad u,v\ge0.}
\]
当 $u+v$ 为定值时，可直接求得最大值；若遇 $\sqrt{A}+\sqrt{B}$ 型最值，也可通过均值不等式 $\sqrt{uv}\le\frac{u+v}{2}$ 或三角换元转化为三角函数有界性。

\textbf{使用场景：}求形如 $\sqrt{ax+b}\pm\sqrt{cx+d}$、$\sqrt{a^{2}-x^{2}}\pm x$ 或含 $\sqrt{x^{2}+a^{2}}$ 的函数最值，常见于高考选填或解答题中的最值步骤。

\textbf{简单例子：}求 $f(x)=\sqrt{x-1}+\sqrt{5-x}$ 的最大值。
解：由 $x-1\ge0,\,5-x\ge0$ 得 $x\in[1,5]$。令 $u=x-1,\;v=5-x$，则 $u+v=4$。
由核心公式，$f(x)\le\sqrt{2(u+v)}=\sqrt{8}=2\sqrt{2}$。
当 $u=v$ 即 $x-1=5-x\implies x=3$ 时取等号，故最大值为 $2\sqrt{2}$。